\documentclass[11pt]{article}
\usepackage{amsmath,amssymb,amsthm,geometry,hyperref}
\geometry{margin=1in}
\newtheorem{theorem}{Theorem}
\begin{document}
\title{An analytical fourth-order anticoherent state on nineteen qubits}
\author{Exact catalogue construction}
\date{3 November 2026}
\maketitle
\begin{theorem}
There exists a pure symmetric nineteen-qubit state of exactly fourth
anticoherence order. The separate \href{exclusion/proof.tex}{exact plane exclusion} then proves $K(19,4)=1$. Together with the independently certified
nonexistence of fifth-order states, the maximum pure-state AC order is four.
\end{theorem}
In the normalized Dicke basis, set
\[
 |\psi\rangle=\frac{\sqrt{194}}{40}|D_{19}^{(0)}\rangle
 +A|D_{19}^{(6)}\rangle+\frac{\sqrt{266}}{160}|D_{19}^{(10)}\rangle
 +\frac{\sqrt{7410}}{160}|D_{19}^{(14)}\rangle+B|D_{19}^{(18)}\rangle,
\]
where
\begin{align*}
 A&=-\frac{9673\sqrt{17290}}{5605600}
       -i\frac{\sqrt{2918881141085}}{2802800},\\
 B&=-\frac{2837\sqrt{323323}}{10210200}
       -i\frac{\sqrt{218332309353158}}{40840800}.
\end{align*}
This is a new exact witness, not an identification of the earlier
seven-component numerical representative from \texttt{ACLiege.jl}.

\section*{Rational quadrature and a phase triangle}
The five probabilities, at indices $S=(0,6,10,14,18)$, are
\[
 p=\left(\frac{97}{800},\frac{1083}{2560},\frac{133}{12800},
         \frac{741}{2560},\frac{399}{2560}\right).
\]
Direct rational substitution gives
\[
 \sum_{r\in S}p_r r^a=\frac1{20}\sum_{r=0}^{19}r^a,
 \qquad 0\le a\le4.
\]
Their common values are $1,19/2,247/2,1805,281333/10$.
Thus the state is normalized and every diagonal moment through degree
four has its isotropic value. No two support indices differ by one,
two or three. Hence all nonzero magnetic shifts below four vanish.

At shift four the only edges are $6\to10$, $10\to14$ and $14\to18$.
Put $\Lambda_r=\sqrt{\prod_{a=1}^4(r+a)(20-r-a)}$ and define the
three positive lengths
\[
 a=\sqrt{p_6p_{10}}\Lambda_6=\frac{1197\sqrt{27170}}{320},\quad
 b=\sqrt{p_{10}p_{14}}\Lambda_{10}=\frac{15561\sqrt{2310}}{1600},\quad
 c=\sqrt{p_{14}p_{18}}\Lambda_{14}=\frac{513\sqrt{1547}}{32}.
\]
Let
\[
 x=\frac{c^2-a^2-b^2}{2ab}=-\frac{9673\sqrt{5187}}{1996995},
 \qquad u=x+i\sqrt{1-x^2},\qquad v=-\frac{au+b}{c}.
\]
Here $-1<x<1$. The cosine-rule identity gives $|u|=|v|=1$
and $au+b+cv=0$. Choose
\[
 c_0=\sqrt{p_0},\quad c_6=\sqrt{p_6}\,\overline u,\quad
 c_{10}=\sqrt{p_{10}},\quad c_{14}=\sqrt{p_{14}},\quad
 c_{18}=\sqrt{p_{18}}v.
\]
These simplify to the radicals stated above. They obey
\[
 \sum_{r=0}^{15}\overline c_r c_{r+4}\Lambda_r=au+b+cv=0.
\]
This and its conjugate are the remaining rank-four off-diagonal
conditions. Together with the diagonal conditions and the absent shifts,
they exhaust every irreducible spin-tensor component of ranks one
through four. Consequently $\rho_4=I_5/5$.

\section*{Independent direct reduction check}
The verifier reconstructs all entries of the reductions from
\[
 (\rho_t)_{ab}=\sum_{\ell=0}^{19-t}c_{a+\ell}\overline c_{b+\ell}
 \sqrt{\frac{\binom ta\binom tb\binom{19-t}{\ell}^{2}}
 {\binom{19}{a+\ell}\binom{19}{b+\ell}}}.
\]
Exact fractions and squarefree radicals give $\rho_t=I_{t+1}/(t+1)$
for $t=1,2,3,4$, while
\[
 \left\|\rho_5-I_6/6\right\|_F^2=\frac{8742251}{10653696000}>0.
\]
Thus the representative has exactly order four. The separate
\href{../n19-t5/proof.tex}{four-copy positive-moment certificate} excludes every
fifth-order pure state, proving global optimality of order four.
The separate \href{exclusion/proof.tex}{Young-space positive-moment certificate} excludes every two-dimensional AC4 subspace, including arbitrary complex bases. Consequently this exact line is also dimension-optimal: $K(19,4)=1$.
\end{document}
